\section*{Responsible AI Statement}

This research demonstrates the application of AI-driven optimization to parallel computing, specifically targeting GraphViz layout algorithms with OpenMP parallelization. We acknowledge both the potential benefits and risks associated with AI-generated code optimization and provide this statement to address broader impacts and ethical considerations.

\textbf{Positive Societal Impacts:} Our AI-driven approach democratizes parallel computing optimization by reducing the expertise barrier for achieving high-performance implementations. This can accelerate scientific computing across diverse domains, from computational biology to climate modeling, enabling researchers without specialized parallel programming knowledge to leverage multi-core architectures effectively. The automated optimization pipeline can significantly reduce development time and improve computational efficiency, leading to energy savings and reduced computational costs in large-scale scientific applications.

\textbf{Potential Risks and Mitigation Strategies:} We recognize several potential risks: (1) \textit{Code correctness concerns} - AI-generated parallel code may introduce subtle race conditions or synchronization errors. We mitigate this through comprehensive validation using ThreadSanitizer, Valgrind, and extensive correctness testing across diverse graph configurations. (2) \textit{Over-reliance on automation} - Researchers may become overly dependent on AI optimization without understanding underlying parallel programming principles. We address this by providing detailed explanations of optimization strategies and maintaining transparency in our AI decision-making process. (3) \textit{Performance regression risks} - Automated optimizations may not always improve performance across all scenarios. Our statistical validation methodology with 95\% confidence intervals and comprehensive benchmarking across diverse workloads helps identify and prevent such regressions.

\textbf{Ethical Considerations:} Our research adheres to the NeurIPS Code of Ethics and emphasizes transparency, reproducibility, and responsible deployment. We provide complete source code, detailed experimental protocols, and comprehensive documentation to enable independent verification and responsible use of our methods. The AI ensemble approach includes multiple validation layers to ensure reliability and reduce the risk of generating incorrect or harmful optimizations.

\textbf{Safe Deployment Practices:} We recommend that practitioners using AI-generated parallel code: (1) conduct thorough testing with representative workloads, (2) validate correctness using appropriate tools, (3) benchmark performance against baseline implementations, and (4) maintain human oversight in production deployments. Our methodology provides a framework for responsible AI-assisted optimization that balances automation benefits with necessary safety measures.
